\section{The Craft}

For roughly a century and a half, from the early fifteenth century to the Counter-Reformation, the best-selling religious book in Europe after the Bible was a short manual on how to die. It circulated in Latin under at least three titles, in a picture-book version with eleven woodcuts, in French, in German, and---in the version this section will use---in English, where it was called \emph{The Book of the Craft of Dying}.\endnote{\emph{The Book of the Craft of Dying, and Other Early English Tracts Concerning Death}, taken from manuscripts and printed books in the British Museum and Bodleian libraries, ``now first done into modern spelling and edited by Frances M.~M. Comper'' (London: Longmans, Green, and Co., 1917); read 2026-07-26 in the Internet Archive's digitization of the 1917 printing. The essay quotes Comper's modernized-spelling text, not a manuscript, and this is a real limitation: Comper transcribed Bodleian MS 423 (which she judged the earliest of three Bodleian witnesses), collated it against Douce MS 322 and Rawlinson MS C.~894, and normalized the spelling throughout. The English tract is itself a translation of a Latin original. The wording quoted here is therefore at two removes from the Latin and one from the manuscript, and no claim is made about the Latin's wording except where Comper reports it.}

Its subject was not death in general. It was the last few hours, and what to say and do in them.

\subsection{Who did not write it}

The tract is very frequently attributed to Jean Gerson, chancellor of the University of Paris, and the attribution does not survive contact with the text.

Comper, who did the collation, sets out the reasons. Gerson certainly compiled a treatise in Latin and French called the \latin{Opusculum tripartitum}, whose third part is \latin{de arte moriendi}. But that book ``is very much shorter than the English version of \emph{The Craft of Dying}, and there is nothing in it which corresponds to the first two chapters of the Craft.'' And then the decisive point: the Craft itself, in its third and fifth chapters, cites the authority of ``the noble'' and ``great clerk, the Chancellor of Paris''---which is to say, it quotes Gerson as a source.\endnote{Comper, ``Note on the Book of the Craft of Dying,'' in \emph{The Book of the Craft of Dying}, pp.~48--51; read 2026-07-26 as above. She adds that the tract ``has been ascribed to Richard Rolle'' as so many English writings of the period were, that ``it may possibly have been translated by him into English, but the author of the older Latin original is unknown,'' and that ``the whole question of the authorship and the various versions of the treatises which are in the catalogues generally included under the title \latin{Ars Moriendi} is one of some difficulty and obscurity.'' She distinguishes at least three distinct books: the Latin treatise (found as \latin{De Arte Moriendi}, \latin{Tractatus de Arte Moriendi}, and \latin{Speculum Artis Moriendi}), the block-books, and the French \emph{L'Art de bien Vivre et bien Mourir}, a different work published by V\'erard at Paris in 1493. This essay reports her findings as those of an identified editor working from the manuscripts, not as its own collation.}

A book does not cite itself as an authority in the third person. So the anonymous Latin author drew on Gerson, and the tradition later paid him the compliment of assuming he \emph{was} Gerson. That is a small point of bibliography, and it is worth making because the \emph{ars moriendi} is regularly presented as the considered teaching of a great churchman, when its actual status is that of an anonymous compilation of enormous popularity---which is a different kind of evidence about what the late medieval Church believed, and in some ways a better one.

\subsection{What is in it}

Six chapters: a commendation of death; the temptations of the dying; the questions to ask them; an instruction with prayers modeled on Christ's five acts on the cross; an instruction to those who will die; and prayers to be said over the dying ``of some man that is about them.''\endnote{\emph{The Craft of Dying}, table of contents and prologue, Comper 1917, pp.~3--4; read as above. The prologue's own summary of the sixth chapter---prayers ``that should be said to them that be a-dying, of some men that be about them''---is quoted in the body.}

The first chapter is a commendation, and its logic is not what a modern reader expects. It begins by conceding Aristotle's point that bodily death is ``most dreadful of all fearful things,'' and immediately outranks it: the death of the soul is worse in the same proportion that the soul outranks the body. Then two psalm verses set against each other---``the death of the sinful man is worst of all deaths'' and ``the death of the good man is ever precious in the sight of God''---and the conclusion that what makes a death good or bad is not its manner: ``what manner of bodily death that ever they die.''\endnote{\emph{The Craft of Dying}, ch.~1, Comper 1917, pp.~5--6; read as above. The tract cites ``the Philosopher\ldots{} in the third book of Ethics'' for the dreadfulness of death, Psalm 33:22 (\latin{Mors peccatorum pessima}) and Psalm 115:15 (\latin{Preciosa est in conspectu Domini mors sanctorum eius}) in Vulgate numbering, Apocalypse 14:13, and Wisdom 4:7. All four were checked in the Douay--Rheims tracked verse text 2026-07-26 and correspond: Psalm 33:22, ``The death of the wicked is very evil''; Psalm 115:15, ``Precious in the sight of the Lord is the death of his saints''; Apocalypse 14:13, ``Blessed are the dead who die in the Lord''; Wisdom 4:7, ``But the just man, if he be prevented with death, shall be in rest.''}

Then comes a passage that a careful reader should stop at, because it is the point where the devotional tradition and the dogmatic one part company. Death, says the tract, quoting ``a wise man,'' is ``nothing else but a going out of prison, and an ending of exile; a discharging of an heavy burden, that is the body; finishing of all infirmities; a scaping of all perils; destroying of all evil things; breaking of all bonds; paying of the debt of natural duty; turning again into his country; and entering into bliss and joy.''\endnote{\emph{The Craft of Dying}, ch.~1, Comper 1917, p.~6; read as above. The tract immediately restricts the commendation---``And this is understood only of good men and the chosen people of God''---and quotes Ecclesiastes 7:2 (``The day of man's death is better than the day of man's birth''), which the Douay renders ``A good name is better than precious ointments: and the day of death than the day of one's birth'' (checked 2026-07-26). The phrase ``paying of the debt of natural duty'' is the English tradition's rendering of the \latin{debitum naturae} from which this essay takes its title.}

Prison; exile; a heavy burden, \emph{that is the body}. That is Plato, not the Council of Vienne. A body which is a burden to be discharged is not a body whose soul is its form \latin{per se et essentialiter}, and the man who has learned from the Catechism that ``spirit and matter, in man, are not two natures united, but rather their union forms a single nature'' will find his devotional inheritance pulling the other way. The pull is not a heresy---the tract is describing a felt relief, not proposing an anthropology---but it is a drift, and it is worth naming because it is the standing drift of the whole genre. Augustine, who could describe the wrenching apart of body and soul as ``jarring horridly on nature,'' is the better guide here than the anonymous author who found it a discharge.

The second chapter is the famous one: five temptations that assail the dying with a violence they have never met before---``such as they never had before in all their life.'' Against faith, the devil works to unsettle belief or to insinuate doubt. Against hope, he raises despair by ``objecting his sins against him.'' Against charity, impatience, and here the tract is medically observant: the worst cases are those who ``die not by nature and course of age---that happeth right seldom, as open experience teacheth men---but die often through an accidental sickness; as a fever'' or an abscess.\endnote{\emph{The Craft of Dying}, ch.~2, Comper 1917, pp.~9--20; read as above. The five are enumerated in the text: faith, despair, impatience, complacence (spiritual pride), and attachment to temporal things. The aside about how rarely men die of old age is the tract's own, and is a small piece of demographic realism about the fifteenth century, not a claim this essay endorses about any other period.} The fourth temptation is the one that shows the author knew his readership: complacence, ``that is spiritual pride, with the which the devil tempteth and beguileth most religious, and devout and perfect men''---the temptation, at the very end, to admire one's own performance. \emph{O how stable art thou in the faith! how strong in hope! how sad in patience! O how many good deeds hast thou done!} The fifth is attachment to ``their wives, their children, their carnal friends, and their worldly riches.''

The third chapter supplies a script. On the authority of ``Saint Anselm the bishop,'' the man at the bedside is to put a series of questions, each answered \emph{Yea}: Art thou glad that thou shalt die in the faith of Christ? Knowest thou well that thou hast not done as thou shouldst have done? Repentest thee thereof? Hast thou full will to amend thee, if thou mightest have full space of life? Believest thou fully that Our Lord Jesu Christ, God's Son, died for thee? Thankest thou Him thereof with all thine heart? Believest thou verily that thou mayest not be saved but by Christ's death and His passion? And then the instruction: ``put all thy trust in His passion and in His death only, having trust in none other thing.''\endnote{\emph{The Craft of Dying}, ch.~3, Comper 1917, pp.~22--23; read as above. The attribution to Anselm is the tract's own and this essay has not verified it against Anselm's works; what is verified is that the tract makes it. The interrogations are a fixed form which the tract then qualifies: ``though these interrogations abovesaid be commonly'' put, they may be adapted.}

What is striking about the script is its economy. Seven questions, each answerable by a monosyllable, designed for a person who may have very little breath and less attention. It is a liturgy for someone who can no longer perform one.

\subsection{The book's quarrel with the sickroom}

The fifth chapter contains the tract's real argument, and it is a complaint.

``It is greatly to be noted, and to be taken heed of, that right seldom that any man---yea among religious and devout men---dispose themselves to death betimes as they ought. For every man weeneth himself to live long, and troweth not that he shall die in short time; and doubtless that cometh of the devil's subtle temptation.''\endnote{\emph{The Craft of Dying}, ch.~5, Comper 1917, p.~32; read as above. ``Weeneth'' is supposes; ``troweth'' is believes.} Then the practical grievance. A ``certain decretal,'' the tract says, ``chargeth straitly every bodily leech that he give no sick man no bodily medicine unto the time that he hath warned and induced him to seek his spiritual leech. But this counsel is now for-slothed almost of all men, and is turned into the contrary; for men seek sooner and busier after medicines for the body than for the soul.''\endnote{\emph{The Craft of Dying}, ch.~5, Comper 1917, pp.~32--33; read as above. The ``certain decretal'' is not identified in the text, and this essay has not traced it; the reference is reported as the tract's own claim. ``For-slothed'' is Comper's spelling; her footnote glosses the related form as ``lost themselves through sloth.''}

And then the sharpest thing in the book, and the reason it is worth reading in an age with better medicine and worse manners about dying: ``as the great Clerk, the Chancellor of Paris saith: Often times by such a vain and a false cheering and comforting, and feigned behoting''---promising---``of bodily heal, and trusting thereupon, men run and fall into certain damnation everlastingly.''\endnote{\emph{The Craft of Dying}, ch.~5, Comper 1917, p.~33; read as above. Comper's footnote glosses ``behoting'' as ``promising.'' The reference to the Chancellor of Paris is Gerson, and is one of the two citations that establish the tract is not his work.} The complaint is that everyone in the room is lying, in the kindest possible way, and that the lying has a cost the liars are not paying.

Two further features of the tract's construction should be noticed, because they explain what kind of book it is. The prologue says it is written for those who lack skill in dying, ``not only to lewd men but also to religious and devout persons''---and Comper's gloss on \emph{lewd men} is \emph{laymen}. And the sixth chapter's prayers are to be said over the dying ``of some man that is about them''; elsewhere the tract directs that the dying be exhorted ``by a priest, or for need by another.''\endnote{\emph{The Craft of Dying}, prologue (Comper 1917, p.~3, with her footnote glossing ``lewd men'' as ``laymen''), the prologue's summary of ch.~6 (p.~4), and the exhortation formula in the appended tract at p.~85 (``to exhort him, by a priest, or for need by another, in the manner as followeth''). Read 2026-07-26 as above. The inference drawn in the body---that the book presupposes situations in which no priest is present---is this essay's reading of these provisions and is offered as such; the essay has not investigated the demographic or ecclesiastical conditions that produced the genre.}

\emph{Or for need by another.} That is a manual written for a world in which the priest might not come. It hands a layman a script, a diagnosis of the five ways the dying mind fails, and a set of questions to ask his neighbour. Whatever else it is, it is a work of lay competence in a domain the modern Church has largely professionalized---and it was the most popular book of its kind in Christendom.

\subsection{What the Church kept}

The genre continued: Caxton printed an abridgment ``translated from the French'' before 1500, block-books circulated in almost every European language, and the Counter-Reformation produced its own successors.\endnote{Comper, ``Note on the Book of the Craft of Dying,'' pp.~48--51, and the abridgment printed at pp.~55--91 under the title ``The Art and Craft to Know Well to Die,'' which Comper attributes to Caxton and describes as translated from a French source she was unable to trace; read 2026-07-26 as above. She notes that the block-books ``generally contain eleven illustrations depicting the five great temptations which beset the soul at death,'' and that the frontispiece of her edition is taken from the British Museum's block-book published at Cologne c.~1450. This essay makes no claim about the later Counter-Reformation successors, which it has not examined.} What is more interesting is how much of the counsel the Catechism still carries, and where it has quietly changed the emphasis.

``The Church encourages us to prepare ourselves for the hour of our death. In the litany of the saints, for instance, she has us pray: `From a sudden and unforeseen death, deliver us, O Lord'; to ask the Mother of God to intercede for us `at the hour of our death' in the Hail Mary; and to entrust ourselves to St.~Joseph, the patron of a happy death.''\endnote{CCC 1014, Holy See English web text, read 2026-07-26. The paragraph then quotes \emph{The Imitation of Christ} I.23.1 (``Every action of yours, every thought, should be those of one who expects to die before the day is out\ldots{} If you aren't fit to face death today, it's very unlikely you will be tomorrow'') and the stanza of Francis of Assisi's Canticle of the Creatures on ``our sister bodily Death.'' The Catechism's own footnotes identify both sources.}

The petition is worth stopping over, because it inverts a modern preference so completely that most people who pray it do not hear what they are asking. \emph{From a sudden and unforeseen death, deliver us.} Ask a hundred people today how they would like to die and the majority will say: quickly, in my sleep, without knowing. The Church's litany asks to be spared exactly that. Not because slow dying is better---the tradition is under no illusions about what a fever or an abscess does---but because a death one can see coming is a death one can do something about, and the entire \emph{ars moriendi} rests on the assumption that there is something to be done.

That assumption is the tradition's real bequest, and it is not obviously false. What has changed is where the work is located. The medieval tract puts the decisive struggle in the last hours and gives most of its pages to them, while conceding in its fifth chapter that men ought to dispose themselves ``betimes.'' The Catechism reverses the proportions: one paragraph on the hour, and an entire moral theology on the life. The reversal is not a rejection. It is what happens when a counsel that began as emergency instruction is absorbed into an account of the whole Christian life, and it is the same movement by which the tract's own author already knew that the best preparation for the deathbed happens years before anyone reaches it.

None of the foregoing is magisterial teaching, and it should not be cited as though it were. The \emph{ars moriendi} is devotional literature of uncertain authorship, of enormous historical influence, and of no doctrinal authority whatever; its five temptations are a psychology, not a revelation; its interrogations are a pastoral form, not a rite. What the Church has authoritatively adopted from the tradition amounts to CCC 1014's single paragraph and the liturgical petitions it quotes.

But the Church did not leave the dying to a devotional literature. It has, for the whole of its history, done specific things at deathbeds, and those things are law.
